% ======================================================================
% Section 4 — The positive core
% ======================================================================

\section{The positive core: classification at the rigid primes and at
\boldmath$p^2$/$pq$}\label{sec:positive}

The program's attacks were launched from a positive base, and that base
survives the negative results intact. This section states it with
precision; the full proofs, the verification records, and the
measurement methodology are the subject of the companion monograph
\cite{prog-artifacts}, and nothing here improves or weakens them. What
matters for this paper's thesis is that the positive core is exactly as
bounded as it is: everything it proves is a closure, classification, or
identity, and its one conditional theorem wears its hypotheses on its
sleeve.

\subsection{The sum-only equality theorem}

The analytic folklore (recorded in a comment on Tao's weblog
\cite{Tao-blog}) is that the optimal time is rational with denominator
the sum or difference of two speeds. The sharpening the program proved
is that \emph{sums alone suffice}.

\begin{theorem}[Sum-only equality; Theorem 1.1 of the companion
monograph]\label{thm:equality}
Let $V=(v_1,\dots,v_n)$ have distinct positive entries with
$\gcd(V)=1$ and $0<M(V)<\tfrac12$. For each unordered pair $\{u,w\}$
let $\tau(u,w)$ denote the optimal loneliness of the $(n-1)$-runner
instance obtained by restricting to the pair-sum lattice of denominator
$N=v_u+v_w$---explicitly (Appendix~\ref{sec:self}),
\[
\tau(u,w)\;=\;\max_{k\in\Z_{N}}\;\min_{p\in[n]\setminus\{w\}}
\wnc{\tfrac{v_pk}{N}}
\]
(runners $u$ and $w$ coincide on the lattice, since
$\wnc{v_uk/N}=\wnc{v_wk/N}$ there). Then
\[
M(V)\;=\;\max_{\{u,w\}}\tau(u,w),
\]
and consequently $\lrc_n$ holds for all $V$ if and only if the finite
statement $(\taust_n)$ does: for every $V$, some pair-sum lattice
carries an optimal value $\ge 1/(n+1)$. The finitization is exact; no
information is lost, and nothing about the classical difficulty is
claimed to become easier.
\end{theorem}

The two ingredients are that the binding runners at a global optimum can
be chosen two-sided, forcing the optimal time onto a \emph{sum} lattice
(difference denominators never carry the optimum), and a mirror lemma:
on that lattice the two binding runners are exact reflections of each
other at every lattice time, so the $n$-runner minimum there equals an
$(n-1)$-runner minimum. The equivalence of Corollary form
($\lrc_n\iff\taust_n$) is the lossless gateway through which every
instance of \S\ref{sec:tm-covering} entered.

\subsection{The covering framework and the rigid zone}

Under the equality theorem, a pair \emph{certifies} an instance unless
its $n-1$ bad sets $B_w=\{k\in\Z_N:(n{+}1)\wn{wk}<N\}$ cover $\Z_N$.
The framework that follows is finite and exact, and its proved pieces
include a size law for the bad sets, a polygonal witness lemma, moat
counting, a coincidence-counting lemma, a lift law, and a class-group
reduction at prime moduli (Lemma 2.2--2.7 of the companion
monograph). Five empirically distinct obstruction classes from the
program's early experiments collapse, in this language, into additive
orders of one covering phenomenon.

\begin{theorem}[Rigid-zone classifications; Theorems 3.1--3.5 of the
companion monograph]\label{thm:rigid}
At rung $n=4$ the unit-capable moduli---those $N$ for which $n-1$ unit
bad sets can cover $\Z_N$---are $7,11,13$; at rung $n=5$ they are
$7,13,17,19,37$ within the verified range (primes $\le150$, composites
$\le111$). At every one of these primes the covering configurations are
classified completely by a short explicit list of translate classes in
the cyclic group: antipodal classes and domino tilings at $7$ and $11$;
perfect matchings at $13$; the classes $\{0,1,2,7\}$, $\{0,2,4,6\}$ at
$17$ and $19$ (together with one further class at $19$); and the single
class $\{0,2,12,15\}$ at $37$.
\end{theorem}

This is the program's central structural finding: \emph{at rigid primes,
the covering obstruction is governed by a small translate-class
structure}, and the classification is what closes $75.6\%$ of the
$3{,}712{,}576$-instance corpus outright. The residual obstruction is
localized (mixed order signatures at composite moduli), and the composite
side carries the structure that the next subsection records.

\subsection{The composite side: group ring, tensor, and the fiber
cascade}

Beyond the unit capable moduli, at $N=p^2$ and $N=pq$, the program built
the composite-side structure that the scaling question lives on. The
group-ring translation expresses the covering count as a cyclotomic
coefficient problem; the $pq$ side decomposes as a tensor problem whose
Lam--Leung-type structure \cite{LamLeung} fixes the admissible
fiber-count signatures; and the foot-count input to the cascade was
discharged by an exact identity verified on $10{,}410$ triples with the
fixed threshold at $p=31$. The higher rungs carry a general-$T$ fiber
machinery, an exact dilate-intersection lemma, a reflection-fiber
invariance, a chain-class dilate theorem, and a renormalization lemma;
the Mirsky--Newman covering-system theorem and the
R\'edei--de Bruijn--Schoenberg structure theorem for vanishing sums of
roots of unity \cite{MirskyNewman,LamLeung} play the classical
supporting roles.

Everything above is proved and machine-verified. What is \emph{not}
proved is the scaling closure itself, and the program consolidated it
into a single conditional statement with a disclosed ledger.

\begin{theorem}[Conditional scaling closure; Theorem 5.1 of the
companion monograph]\label{thm:SC}
Assume the distinct-family volume law H1g and the dilate-sparsity pair
H2m/H2r. Then the fiber cascade closes the covering obstruction at the
frontier cells $(1K,T{-}3U)$ of every rung $T\ge8$ at $N=p^2$, within a
depth that is $T$-flat in practice and never exceeds $T+3$ in the worst
case. The reduction chain and every other input are proved or
machine-validated; each constant's provenance---derived, measured, or
fitted, with the data it is fit to---is disclosed in the six-entry
hypothesis ledger, and the two obligations are reduced to named objects:
the strand census (for H1g) and the sharp pair-line constants (for
H2m).
\end{theorem}

We do not restate the ledger here. Its relevance to this paper is its
shape: the proved side of the ledger is all closures and identities; the
open side is exactly the two growth-in-$k$ margins whose type-mismatch
career occupies \S\ref{sec:tm}. Theorem~\ref{thm:SC} is what the positive
core amounts to when stated at its honest strength, and it is why the
program's claim to a positive result is deliberately the bounded one:
\emph{the obstruction is completely classified where classification
suffices, and the residual is a named margin of a type no turn of the
program ever proved}.
